\section{A fixed input order with exponentially many affine mistakes}

We give the complete affine construction used by the network theorem.
Keep the target $f_k$ from Theorem~\ref{thm:main}, and abbreviate the two
encoded integers by $X$ and $Y$. The calibration examples all have
selector $z=-1$.
\begin{enumerate}
\item Take diagonal inputs $x=y$ equal to the all-positive vector,
each of its $k$ one-coordinate flips, and the negatives of these vectors.
Every label is positive.
\item For each $i\in\{0,\ldots,k-1\}$, take two negative examples with
$x_i=-1,y_i=+1$ and $x_j=+1,y_j=-1$ for $j<i$. For $j>i$, set the equal
bits $x_j=y_j$ all positive in one example and all negative in the other.
Their highest differing bit is $i$, so $X<Y$.
\end{enumerate}
Remove duplicate examples. The resulting calibration has at most $4k+2$
inputs. Let $q_t$ denote the input $z=+1$, $Y=0$, $X=t$; its target label
is negative for every $0\le t\le2^k-1$.

\begin{lemma}\label{lem:selector}
For $k\le t\le2^k-1$, the unique minimum-norm affine unit-margin
coefficient vector fitting the calibration and the negative example
$q_t$ is the vector of the affine score
\begin{equation}\label{eq:selector-score}
 a_t\cdot\bar x=-t-(t+1)z+\sum_{i=0}^{k-1}2^i(x_i-y_i).
\end{equation}
It fits all earlier $q_j$ with $j\le t$, and its absolute score is at
least one on the entire cube. Its squared norm is
\begin{equation}\label{eq:Qt}
 Q_t=t^2+(t+1)^2+2\sum_{i=0}^{k-1}4^i.
\end{equation}
For $t<2^k-1$, it predicts $q_{t+1}$ positively, and adding the correct
negative label raises the convex affine minimum by
\begin{equation}\label{eq:Qgap}
 Q_{t+1}-Q_t=4(t+1)>0.
\end{equation}
\end{lemma}
\begin{proof}
Write an arbitrary affine score as
\[
 b+cz+\sum_i a_i x_i+\sum_i h_i y_i,
\]
and set
\[
 u=b-c,\quad v=b+c,\quad
 w_i=(a_i-h_i)/2,\quad s_i=a_i+h_i,\quad U=u-1.
\]
The squared coefficient norm is
\begin{equation}\label{eq:affine-norm}
 \frac{u^2+v^2}{2}+2\sum_iw_i^2+\frac12\sum_i s_i^2.
\end{equation}
The paired positive diagonal constraints give
\[
 U\ge0,\qquad \left|\sum_j s_j\right|\le U,\qquad
 \left|\sum_j s_j-2s_i\right|\le U.
\]
Subtracting the expressions inside the last two absolute values shows
$|s_i|\le U$, hence $\sum_i|s_i|\le kU$.
Average the two negative carry constraints at coordinate $i$. The upper
equal-bit terms cancel and leave
\[
 u-2w_i+2\sum_{j<i}w_j\le-1.
\]
Induction yields
\begin{equation}\label{eq:binary-lower}
 w_i\ge\sum_{j<i}w_j+\frac{u+1}{2}
       \ge2^i(1+U/2).
\end{equation}
At $q_t$, the non-selector part of the input contribution is
$2\sum_{\operatorname{bit}_i(t)=1}w_i-
\sum_{\operatorname{bit}_i(t)=0}s_i$. Its negative unit-margin constraint
therefore implies
\[
 v\le-1-2\sum_{\operatorname{bit}_i(t)=1}w_i
            +\sum_{\operatorname{bit}_i(t)=0}s_i
   \le-(2t+1)+(k-t)U\le-(2t+1),
\]
since $t\ge k$. Combining this with $u\ge1$ and
\eqref{eq:binary-lower} in \eqref{eq:affine-norm} gives
\[
 \|a\|^2\ge\frac{1+(2t+1)^2}{2}+2\sum_i4^i=Q_t.
\]
Equality requires $u=1$, $s_i=0$, $w_i=2^i$ and $v=-(2t+1)$,
which is exactly \eqref{eq:selector-score}.

Conversely the displayed score equals $1+2(X-Y)$ on $z=-1$, and
$-(2t+1)+2(X-Y)$ on $z=+1$. It fits every calibration label, and its score
at $q_j$ is $2j-2t-1\le-1$ for $j\le t$. Both branch scores are odd
integers on the full cube, so the absolute margin is at least one.
At $q_{t+1}$ the score is $+1$. Applying the same minimum calculation
to the augmented history gives $Q_{t+1}$, and subtraction gives
\eqref{eq:Qgap}.
\end{proof}

The input permutation starts with the calibration, then $q_k$, then
$q_{k+1},\ldots,q_{2^k-1}$, and finally every remaining input in
lexicographic order. These blocks are disjoint. Its calibration prefix
has length at most $4k+3$, and there are
$T_k=2^k-k-1\ge2^{k-1}$ designated queries. Just before each one, the
wrong affine witness has the exact convex margin minimum for all data
seen so far.

The target has a linear-size Boolean circuit. With $0/1$ bit notation,
comparison of the lowest $i+1$ bits obeys
\[
 G_{-1}=1,\qquad
 G_i=(X_i\wedge\neg Y_i)\ \vee\
       ((X_i\leftrightarrow Y_i)\wedge G_{i-1}).
\]
Each step uses a constant number of bounded-fan-in gates; conjoin the
last result with $z=-1$. Every input bit is essential: for $x_i$ use
$Y=2^i$ and all other $X$ bits zero; for $y_i$ use $X=0$ and all other
$Y$ bits zero; for $z$ take $X=Y$. A fan-in-two circuit with $d$ essential
inputs has at least $d-1$ gates, by counting edges in the connected
ancestry of its output. This proves the matching $\Theta(k)$ bound.
